\subsection{Разложение гильбертова пространства \texorpdfstring{\(H=\overline{\operatorname{Im}(A-\lambda I)}\oplus\ker(A-\lambda I)\)}{H = cl Im(A-lambda I) plus Ker(A-lambda I)}, где \texorpdfstring{\(A\)}{A} --- самосопряженный оператор}

\begin{definition}[Оператор \(A_\lambda\)]
	Пусть \(H\) --- гильбертово пространство над \(\mathbb{C}\),
	\[
	A\in\mathcal{L}(H),
	\qquad
	\lambda\in\mathbb{C}.
	\]
	Обозначим
	\[
	A_\lambda=A-\lambda I.
	\]
\end{definition}

\begin{theorem}
	Пусть \(H\) --- гильбертово пространство над \(\mathbb{C}\), и
	\[
	A\in\mathcal{L}(H)
	\]
	--- самосопряженный оператор. Тогда для любого \(\lambda\in\mathbb{C}\)
	\[
	H=\overline{\operatorname{Im}A_\lambda}\oplus_{\perp}\ker A_\lambda.
	\]
	То есть
	\[
	H
	=
	\overline{\operatorname{Im}(A-\lambda I)}
	\oplus_{\perp}
	\ker(A-\lambda I).
	\]
\end{theorem}

\begin{proof}
	Будем использовать разложение из билета 10.2:
	\[
	H=\overline{\operatorname{Im}B}\oplus_{\perp}\ker B^*
	\]
	для любого
	\[
	B\in\mathcal{L}(H).
	\]
	Применим его к оператору
	\[
	B=A_\lambda=A-\lambda I.
	\]
	Тогда
	\[
	H=\overline{\operatorname{Im}A_\lambda}\oplus_{\perp}\ker A_\lambda^*.
	\]
	Остаётся понять, почему в данном случае
	\[
	\ker A_\lambda^*=\ker A_\lambda.
	\]
	
	\textbf{1. Случай \(\lambda\in\mathbb{R}\).}
	
	Так как \(A\) самосопряжен, то
	\[
	A^*=A.
	\]
	Кроме того,
	\[
	(\lambda I)^*=\overline{\lambda}I.
	\]
	Поэтому
	\[
	A_\lambda^*
	=
	(A-\lambda I)^*
	=
	A^*-\overline{\lambda}I
	=
	A-\overline{\lambda}I.
	\]
	Если \(\lambda\in\mathbb{R}\), то
	\[
	\overline{\lambda}=\lambda.
	\]
	Значит,
	\[
	A_\lambda^*=A-\lambda I=A_\lambda.
	\]
	Следовательно,
	\[
	\ker A_\lambda^*=\ker A_\lambda.
	\]
	Подставляя это в разложение из 10.2, получаем
	\[
	H=\overline{\operatorname{Im}A_\lambda}\oplus_{\perp}\ker A_\lambda.
	\]
	
	\textbf{2. Случай \(\lambda\notin\mathbb{R}\).}
	
	По формуле выше
	\[
	A_\lambda^*=A-\overline{\lambda}I.
	\]
	Так как \(A\) самосопряжен, по билету 11.1 все собственные значения \(A\) вещественны.
	
	Но
	\[
	\lambda\notin\mathbb{R},
	\qquad
	\overline{\lambda}\notin\mathbb{R}.
	\]
	Значит, ни \(\lambda\), ни \(\overline{\lambda}\) не являются собственными значениями оператора \(A\).
	Следовательно,
	\[
	\ker(A-\lambda I)=\{0\},
	\]
	и
	\[
	\ker(A-\overline{\lambda}I)=\{0\}.
	\]
	То есть
	\[
	\ker A_\lambda=\{0\},
	\qquad
	\ker A_\lambda^*=\{0\}.
	\]
	
	Из разложения из 10.2 имеем
	\[
	H=\overline{\operatorname{Im}A_\lambda}\oplus_{\perp}\ker A_\lambda^*.
	\]
	Но \(\ker A_\lambda^*=\{0\}\), значит
	\[
	H=\overline{\operatorname{Im}A_\lambda}.
	\]
	Так как одновременно \(\ker A_\lambda=\{0\}\), можно записать
	\[
	H=\overline{\operatorname{Im}A_\lambda}\oplus_{\perp}\ker A_\lambda.
	\]
\end{proof}